% Chapter 7 — Conclusion and Future Scope

\chapter{Conclusion and Future Scope}
\label{Chapter7}

This chapter concludes the report by evaluating the system's deliverables against the research objectives articulated in Chapter \ref{Chapter1}, summarizing key scientific contributions, addressing clinical and technical limitations, and charting a realistic future roadmap for development and clinical deployment.

%----------------------------------------------------------------------------------------

\section{Achievements Against Research Objectives}

Table \ref{tab:objectives_achieved} provides a point-by-point audit of the seven research objectives formulated in Section \ref{Chapter1}.3, detailing the technical deliverables and empirical outcomes achieved.

\begin{table}[h]
\caption[Achievements Against Research Objectives]{Audit of Achievements Against Research Objectives (O1--O7).}
\label{tab:objectives_achieved}
\centering
\scriptsize
\begin{tabular}{p{3.2cm}p{5.5cm}p{4cm}c}
\toprule
\tabhead{Objective} & \tabhead{Technical Deliverable} & \tabhead{Empirical Outcome} & \tabhead{Status} \\
\midrule
\textbf{O1: Synthetic Dataset Construction} & \texttt{case\_generator.py}, \texttt{bio\_tagger.py}, \texttt{splitter.py} producing 19,732 cases across 4,933 MIND rubrics. & 99,854 entities; 0 offset drift (100\% alignment); 86.9\% lexical diversity in pilot. & \textbf{ACHIEVED} \\
\midrule
\textbf{O2: Multi-Encoder NER} & Fine-tuned BioLinkBERT, PubMedBERT, Bio\_ClinicalBERT, and DeBERTa-v3 on 15 BIO token labels. & BioLinkBERT achieved top token BIO F1 of 34.86\%; PubMedBERT achieved top LOC F1 (33.80\%). & \textbf{ACHIEVED} \\
\midrule
\textbf{O3: Ensemble \& Hybrid Pipeline} & \texttt{ensemble.py} (IoU span clustering) and \texttt{hybrid\_pipeline.py} with LLaMA 3.1 LLM arbitration. & Resolves multi-model span conflicts; 100\% valid JSON profile extraction. & \textbf{ACHIEVED} \\
\midrule
\textbf{O4: LLM QLoRA Integration} & 4-bit NF4 QLoRA fine-tuning of BioMistral 7B and Qwen 2.5 7B via BitsAndBytes. & BioMistral 7B achieved 47.59\% Micro F1, 100\% schema validity, 99.73\% grounding. & \textbf{ACHIEVED} \\
\midrule
\textbf{O5: Dense Semantic Retrieval} & ChromaDB HNSW cosine index over all 74,513 Kent rubrics using \texttt{all-MiniLM-L6-v2}. & 74,513 rubrics indexed in 931s; Top-5 query latency $<$18ms; Recall@20 = 41.56\%. & \textbf{ACHIEVED} \\
\midrule
\textbf{O6: Clinical Workstation} & Streamlit workstation with 4-step stepper, Totality Matrix, SQLite Case Store, ReportLab PDF. & 195 automated tests passing; multi-page PDF generation in $\sim$420ms; zero data egress. & \textbf{ACHIEVED} \\
\midrule
\textbf{O7: Empirical Evaluation} & Systematic comparative benchmark on held-out test set ($N = 1,973$ cases). & Comprehensive master matrix across 10 metrics; documented complementary strengths. & \textbf{ACHIEVED} \\
\bottomrule
\end{tabular}
\end{table}

Every objective established at project inception was fully realized and validated with empirical evidence.

%----------------------------------------------------------------------------------------

\section{Summary of Scientific and Technical Contributions}

Kent-AI makes the following contributions to biomedical natural language processing and clinical informatics:

\begin{enumerate}
  \item \textbf{First Homeopathic Clinical NER Dataset}: The creation of an open-access corpus of 19,732 clinically annotated patient narratives with 99,854 entity spans across 15 BIO classes, establishing the first benchmark corpus for homeopathic NLP.
  \item \textbf{Subword Alignment Repair Algorithm}: A robust alignment technique resolving character offset drift between LLM-generated span strings and transformer token boundaries, ensuring zero annotation corruption.
  \item \textbf{Cross-Architecture Benchmark}: The first rigorous empirical comparison of four biomedical transformer encoders and two QLoRA large language models on classical repertorial extraction, demonstrating that PubMed literature pre-training substantially outperforms MIMIC-III ICU pre-training for descriptive symptom language.
  \item \textbf{Generative vs. Discriminative Discovery}: Demonstrating that instruction-tuned 7B LLMs achieve +12.73\% higher F1 than 110M-parameter discriminative encoders on low-resource clinical schemas, while encoders provide 100$\times$ faster inference latency.
  \item \textbf{Sub-18ms Dense Repertory Retrieval}: Demonstration that dense vector embeddings over 74,513 rubrics bridge the lexical mismatch between modern patient colloquialisms and 19th-century taxonomic entries.
  \item \textbf{Production-Grade Workstation}: An open-source, offline clinical workstation with 195 passing automated tests, providing an end-to-end framework from conversational intake to PDF documentation.
\end{enumerate}

%----------------------------------------------------------------------------------------

\section{Limitations}

Despite its achievements, the system has several limitations:

\begin{enumerate}
  \item \textbf{Repertory Scope}: While the dense vector index contains all 74,513 rubrics across all 37 chapters, the fine-tuned NER models and synthetic training dataset currently focus on Chapter 1 (MIND, 4,933 rubrics). Multi-chapter physical symptoms are supported via dense retrieval, but entity classification is optimized for mental/emotional cases.
  \item \textbf{Synthetic Training Distribution}: All 19,732 training cases were generated synthetically by LLaMA 3 8B. While validated by high lexical diversity (86.9\%) and zero offset drift, synthetic cases may exhibit stylistic uniformities not present in real clinical practice.
  \item \textbf{Token F1 Ceiling}: The highest encoder BIO F1 of 34.86\% reflects the difficulty of strict exact-boundary token matching on descriptive colloquial language. While BioMistral 7B reaches 47.59\%, significant headroom remains for future model refinement.
  \item \textbf{Hardware Requirements for LLM Mode}: Running BioMistral 7B locally via Ollama requires $\geq$8 GB GPU VRAM. In resource-constrained clinics, the system must fall back to the lighter encoder ensemble or deterministic rule-based dialogue modes.
  \item \textbf{Absence of Clinical Trial Validation}: Kent-AI has been evaluated on held-out synthetic test sets and verified by simulated consultations; it has not yet undergone prospective clinical trials with patient cohorts.
\end{enumerate}

%----------------------------------------------------------------------------------------

\section{Future Scope and Realistic Extensions}

\subsection{Near-Term Extensions (3 to 6 Months)}
\begin{itemize}
  \item \textbf{Dual-Path Hybrid Retrieval}: Implement a hybrid fusion engine combining entity-constrained retrieval (Path A) with full-narrative holistic dense retrieval (Path B), leveraging the superior Recall@20 (62.95\%) observed in DeBERTa-v3 experiments.
  \item \textbf{Multi-Chapter Training Expansion}: Extend synthetic case generation across the remaining 36 chapters (74,513 rubrics) to create a comprehensive 100,000+ case multi-chapter corpus.
  \item \textbf{Production BioMistral Serving}: Deploy the fine-tuned BioMistral 7B adapter as the primary inference engine within the Two-Stage Hybrid Pipeline, replacing the four-encoder ensemble for higher extraction F1.
\end{itemize}

\subsection{Medium-Term Extensions (6 to 18 Months)}
\begin{itemize}
  \item \textbf{Multi-Lingual Intake}: Extend the intake FSM and embedding models to support Hindi, Bengali, and regional Indian languages using multi-lingual sentence transformers (\texttt{paraphrase-multilingual-MiniLM-L12-v2}), enabling accessibility across rural healthcare centers.
  \item \textbf{Prospective Clinical Pilot Study}: Conduct an observational clinical study with licensed homeopathic physicians ($N=50$ patients) to evaluate diagnostic concordances between Kent-AI recommendations and expert human prescriptions.
  \item \textbf{Peer-Reviewed Academic Publication}: Prepare and submit the empirical benchmark findings to the \textit{Journal of Biomedical Informatics} (JBI) or \textit{IEEE Journal of Biomedical and Health Informatics} (JBHI).
\end{itemize}

\subsection{Long-Term Vision}
The long-term vision is an \textbf{open-source, privacy-preserving clinical intelligence platform} that integrates voice-based patient consultation transcription, dense repertorial totality analysis, differential Materia Medica synthesis, and longitudinal patient outcome tracking, democratizing expert homeopathic knowledge while preserving physician oversight.

%----------------------------------------------------------------------------------------

\section{Concluding Assessment}

Kent-AI demonstrates that the union of modern biomedical Natural Language Processing, dense semantic retrieval, and large language model reasoning can effectively modernize classical homeopathic repertorization. By automating the labour-intensive translation of unstructured colloquial patient narratives into structured, grade-weighted rubric totalities in under 20 milliseconds, the system resolves the three historical bottlenecks of consultation duration, cognitive recall limits, and vocabulary mismatch. 

The complete codebase, configurations, model adapters, and evaluation suites are publicly available under the permissive MIT license at \url{https://github.com/Krishna200608/kent-ai}, providing a reproducible foundation for clinical decision support in complementary medicine.
